% ===================================================================
% Supplementary Figures --- v6 diffusion-based inverse-RG model
% Paste into your paper where the appendix/supplementary section goes.
% Assumes: \usepackage{graphicx}, \usepackage{amsmath}, \usepackage{amssymb}
% Images expected at figures/01_overview.png ... figures/26_instanton_distributions.png
% ===================================================================

\section{Supplementary Figures --- v6 Model}
\label{sec:appendix}

This appendix documents the current (v6) checkpoint only. It focuses on the
two questions the project exists to answer: (1) does the diffusion-generated
ensemble match the true physics of 2D compact U(1) lattice gauge theory, at
couplings inside and outside what the model trained on, and (2) how does
producing that ensemble compare, in cost, to standard HMC --- both in how
fast a usable configuration appears, and in what standard HMC's own
``equilibrium'' ensemble actually contains at the same coupling.

\subsection{Methodology}
\label{sec:appendix-methods}

\paragraph{The model.} A score network is trained to reverse a noising
process applied to wrapped link angles $\theta \in (-\pi, \pi]$, conditioned
on a \emph{coarse} (2$\times$2-blocked) version of the same field.
Generation runs that reverse process backward from noise, which is exactly
one inverse renormalization group step: it lifts a coarse-lattice
configuration, produced by ordinary HMC at a coarse coupling $\beta_c$, up
to a fine-lattice configuration at a target coupling $\beta_f$ (tree-level
relation $\beta_c = \beta_f / 4$). Two mechanisms make this work outside the
training set:

\begin{itemize}
  \item \textbf{Exact-score blend at small noise.} As the reverse process
    approaches zero noise, the learned score is blended toward the
    \emph{analytically exact} score of the (Gaussian-smeared) Wilson
    action. Near $\sigma = 0$ this makes the last steps of generation
    literal Langevin dynamics under the true action --- correct by
    construction, regardless of what the network learned from data. This is
    the main reason the model can be asked for couplings the training set
    never covered: the fine-detail correctness at small noise comes from
    the known physics of the action, not from interpolating training
    examples.
  \item \textbf{Structural topology transport.} The coarse ensemble's
    topological charge $Q_c$ is known exactly (it is an integer, computed
    directly from the coarse configuration) and is deterministically
    projected onto the generated fine configuration, rather than left to
    emerge from the sampling dynamics. This matters because standard HMC
    \emph{cannot} change $Q$ reliably at large $\beta$
    (Section~\ref{sec:appendix-therm}) --- by not relying on stochastic
    tunneling for its topological content, the pipeline sidesteps HMC's
    worst failure mode entirely.
\end{itemize}

\paragraph{Training coverage, precisely.} Training couplings are drawn
continuously, log-uniform in $\beta$ (\texttt{data.random\_rungs}, see
\texttt{diffusion.utils.expand\_rungs}), not from a fixed discrete list: 64
rungs at $L=16$ spanning $\beta \in [1.0, 57.2]$, 12 at $L=32$ spanning
$\beta \in [2.2, 52.6]$, 6 at $L=8$ spanning $\beta \in [1.4, 2.5]$, plus 4
fixed high-statistics anchors at $L=16$ ($\beta = 14.15, 25, 40, 55.02$).
Inside $\beta \approx [1, 60]$, any requested coupling sits close to some
training draw --- this is \textbf{in-sample} generalization (interpolation).
$\beta > 60$ is where training coverage stops entirely --- this is
\textbf{out-of-sample} generalization (extrapolation). Lattice sizes $L=64$
and $L=128$ were never trained on at any coupling --- a separate, orthogonal
out-of-sample axis (volume).

\paragraph{How a figure is derived.} Every distribution-match figure is
produced by \texttt{diffusion.validate.report.validate\_ensemble}: it
measures plaquette angles, Wilson loops (areas 1 to 144), and topological
charge on the generated batch, compares each to the \textbf{exact}
closed-form prediction (2D compact U(1) admits an exact character-expansion
solution for the plaquette distribution, Wilson loop expectations, and the
finite-volume $P(Q)$ --- \texttt{diffusion.lgt.exact} --- no Monte Carlo
involved on the ``exact'' side at all), and reports a $z$-score
$z = (\text{measured} - \text{exact})\,/\,(\text{statistical error of the
measurement})$. $|z| \lesssim 2$ is a pass. A reference HMC ensemble is also
shown for visual comparison; \emph{which} HMC reference matters and is
stated per figure --- an unbiased reference with instanton ($Q$-hop) updates
for ``does the model match the physics'' figures
(Section~\ref{sec:appendix-gen}), versus the case's own plain HMC with no
topological updates for ``what does standard practice actually produce''
figures (Section~\ref{sec:appendix-therm}), since only the latter can
honestly demonstrate a failure mode of standard HMC. \textbf{One exception:}
Section~\ref{sec:appendix-instanton} is the one place in this appendix where
the instanton move is evaluated as a method in its own right rather than
only as a tool for building unbiased references --- there, ``HMC'' and
``instanton HMC'' are both pure \texttt{diffusion.lgt.hmc.BatchedHMC}
output, with no diffusion model involved at all.

% ===================================================================
\subsection{Generalization: does the generated ensemble match the exact physics?}
\label{sec:appendix-gen}
% ===================================================================

Five cases, chosen to cover low/mid/high coupling and both in-sample and
out-of-sample regimes, plus the orthogonal volume axis.

\begin{figure}[htbp]
  \centering
  \includegraphics[width=0.95\textwidth]{figures/01_overview.png}
  \caption{\textbf{Campaign overview --- every case in one picture.} Three
    panels against target coupling $\beta_f$ (shared log axis), shaded band
    = training range $[1, 60]$. \emph{Left:} $\langle Q^2\rangle$ $z$-score
    of the final pipeline output (after topology enforcement and a short
    rethermalization) against exact. \emph{Center:} the raw diffusion
    sampler's own topology-transport quality, before enforcement is
    applied --- how close a bare sample lands to its coarse parent's
    charge. \emph{Right:} raw-seed thermalization cost in HMC trajectories
    (Section~\ref{sec:appendix-therm}). \textbf{Significance:} the left
    panel is the single most important result in the appendix --- $z$ stays
    inside roughly $\pm 2$ continuously from $\beta_f \approx 1$ to
    $\beta_f \approx 870$, with no discontinuity at the training-range
    boundary (shaded band edge). That absence of a jump is the direct
    visual evidence that the exact-score mechanism
    (Section~\ref{sec:appendix-methods}) is doing real work: nothing
    distinguishes interpolation from extrapolation in this panel, which is
    the point.}
  \label{fig:overview}
\end{figure}

\subsubsection{Low \texorpdfstring{$\beta$}{beta}, in-sample}

\begin{figure}[htbp]
  \centering
  \includegraphics[width=0.95\textwidth]{figures/02_gen_low.png}
  \caption{\textbf{$\beta_f = 1.49$, $L=32$ (case A\_bc0.25).} Four panels:
    per-configuration plaquette-angle histogram against the exact
    infinite-volume density; the topological charge histogram $P(Q)$
    against the exact finite-volume distribution; the Wilson-loop area law
    $-\log\langle W(A)\rangle$ against the exact string tension; the
    connected plaquette-plaquette correlator (expected to be pure noise
    here --- the theory is nearly free at weak coupling, so genuine
    correlations are exponentially small). $z$-scores: plaquette $+0.51$,
    $W(2\times2)$ $+1.01$, $\langle Q^2\rangle$ $-1.80$ --- all comfortably
    passing. \textbf{Physics:} at weak coupling the Boltzmann weight
    $e^{\beta\cos\theta}$ is broad, fluctuations are large, and both
    standard HMC and the diffusion model have an easy time --- this case
    exists to confirm the comparison isn't rigged by only showing hard
    cases. \textbf{Significance:} establishes the honest floor. The model
    has no structural advantage here and doesn't need one; showing it still
    matches exactly is the baseline the harder cases build on.}
  \label{fig:gen-low}
\end{figure}

\subsubsection{Mid \texorpdfstring{$\beta$}{beta}, in-sample (at the edge of the trained anchors)}

\begin{figure}[htbp]
  \centering
  \includegraphics[width=0.95\textwidth]{figures/03_gen_mid.png}
  \caption{\textbf{$\beta_f = 55.0$, $L=32$ (case D\_bc14.1464).} Same
    four-panel layout as Figure~\ref{fig:gen-low}. $z$-scores: plaquette
    $-0.50$, $W(2\times2)$ $+0.97$, $\langle Q^2\rangle$ $-0.40$ ($P(Q)$
    $\chi^2$ $p = 0.34$) --- all passing. \textbf{Physics:} $\beta_f = 55$
    sits right at one of the four fixed high-statistics training anchors
    ($\beta = 55.0237$) --- the topological susceptibility here is already
    small ($\langle Q^2\rangle \approx 0.47$ at this volume), so correctly
    reproducing the \emph{shape} of a narrow $P(Q)$ matters more than at low
    $\beta$. \textbf{Significance:} this is the coupling where standard HMC
    starts to struggle badly (Section~\ref{sec:appendix-therm} shows its
    hot-start chain is completely frozen here) --- confirming the model
    still matches the exact distribution at exactly the point where the
    standard method's own ensemble becomes unreliable is the first half of
    the paper's central claim.}
  \label{fig:gen-mid}
\end{figure}

\subsubsection{High \texorpdfstring{$\beta$}{beta}, out-of-sample + volume extrapolation (flagship)}

\begin{figure}[htbp]
  \centering
  \includegraphics[width=0.95\textwidth]{figures/04_gen_high.png}
  \caption{\textbf{$\beta_f = 218.6$, $L=64$ (case F\_L64\_bc55.0237) ---
    the flagship case.} $z$-scores (mean over 2 independent seeds):
    plaquette $-0.51$, $W(2\times2)$ $-0.98$, $\langle Q^2\rangle$ $+0.90$
    ($P(Q)$ $\chi^2$ $p = 0.74$). \textbf{Physics:} two extrapolation
    stresses at once --- $\beta_f$ is $3.6\times$ the training maximum
    (genuinely out-of-sample), and $L=64$ is $4\times$ the largest lattice
    trained on ($16\times$ the volume). \textbf{Significance:} this is the
    single figure that most directly demonstrates the exact-score mechanism
    working as designed: nothing in this coupling/volume combination was in
    the training set, yet the closed-form exact predictions are matched to
    the same tolerance as the in-sample cases above. This is only possible
    because correctness here comes from the blended analytic score at small
    noise (Section~\ref{sec:appendix-methods}), not from the network having
    memorized nearby examples.}
  \label{fig:gen-high}
\end{figure}

\subsubsection{Deepest coupling extrapolation}

\begin{figure}[htbp]
  \centering
  \includegraphics[width=0.95\textwidth]{figures/05_gen_extreme.png}
  \caption{\textbf{$\beta_f = 872.8$, $L=32$ (case F\_L32\_bc218.58) ---
    15$\times$ the training maximum.} $z$-scores: plaquette $-0.41$,
    $W(2\times2)$ $-1.09$. $\langle Q^2\rangle$ is not a meaningful
    $z$-score here (exact $\langle Q^2\rangle \approx 9\times10^{-4}$ at
    this coupling/volume --- both generated and exact are consistent with
    exactly zero fluctuation, which is what the ``$z=\infty$'' appearing in
    the raw tables at this coupling reflects: a zero-variance coincidence,
    not a defect). \textbf{Physics:} at this coupling the theory is almost
    frozen even in principle --- the true plaquette-angle distribution and
    $P(Q)$ have collapsed to a narrow spike, and the test is whether the
    model's \emph{width}, not just its mean, still tracks the exact
    prediction at a coupling scale nearly two orders of magnitude past
    anything trained. \textbf{Significance:} the outer edge of what this
    checkpoint was pushed to demonstrate. That the plaquette-angle
    density's width still matches the exact curve to visual precision here
    is the strongest single data point for ``the mechanism generalizes,''
    precisely because $15\times$ is far enough that no plausible
    interpolation argument applies.}
  \label{fig:gen-extreme}
\end{figure}

\subsubsection{Volume extrapolation at an in-sample coupling}

\begin{figure}[htbp]
  \centering
  \includegraphics[width=0.95\textwidth]{figures/06_gen_volume.png}
  \caption{\textbf{$\beta_f = 14.15$, $L=128$ (case C\_L128) --- the
    largest lattice tested, $8\times$ the largest lattice trained on.}
    $z$-scores: plaquette $+0.97$, $W(2\times2)$ $+1.11$, $\langle
    Q^2\rangle$ $-1.24$. One Wilson loop size is flagged and excluded from
    the area-law panel because its exact expectation value falls below
    $3\sigma$ of measurement noise at this volume (labeled on the panel ---
    a statistical-power limit of the test, not a model failure).
    \textbf{Physics:} this case isolates the volume axis from the coupling
    axis --- $\beta_f = 14.15$ is solidly in-sample, only $L$ is pushed
    beyond training ($L=16/32$ trained $\to$ $L=128$ tested). Finite-volume
    corrections to Wilson loops and $P(Q)$ genuinely change with $L$, so
    this directly tests whether the model has learned volume-dependent
    physics or just a fixed-size lookup. \textbf{Significance:} shows the
    two extrapolation axes (coupling in Figures~\ref{fig:gen-high}--\ref{fig:gen-extreme},
    volume here) are independent successes, not one result restated --- the
    model generalizes in both directions it was asked to.}
  \label{fig:gen-volume}
\end{figure}

\clearpage
% ===================================================================
\subsection{Thermalization: cost vs. standard HMC, and why standard HMC struggles}
\label{sec:appendix-therm}
% ===================================================================

\paragraph{The physics of why standard HMC fails at large $\beta$.} Two
distinct mechanisms, both worsening with $\beta$:

\begin{enumerate}
  \item \textbf{Critical slowing down} --- the integrated autocorrelation
    time of ordinary (non-topological) observables like the plaquette or
    small Wilson loops grows with $\beta$, so a chain needs more
    trajectories between independent samples even once it has reached
    equilibrium.
  \item \textbf{Topological freezing} --- the topological charge $Q$
    changes only via a global rearrangement of the field (an
    instanton-like tunneling event). The rate of these events is
    suppressed roughly as $e^{-2\beta}$ with the action barrier between
    sectors, so above a fairly modest coupling a standard HMC chain's $Q$
    simply stops moving for any practical trajectory budget. This is
    qualitatively different from (1): it isn't slow mixing, it's zero
    mixing --- the chain is stuck in whatever sector it happened to be in,
    forever, as far as any feasible run is concerned.
\end{enumerate}

Standard HMC also has \textbf{two distinct costs}, and figures below are
careful to label which one is meant: \textbf{burn-in} (the one-time cost of
reaching equilibrium from a fresh hot or cold start) and the \textbf{sampling
interval} $2\tau_{\mathrm{int}}$ (once \emph{already} equilibrated, the
trajectories a chain must run and discard before its \emph{next}
configuration is a statistically independent draw --- a different,
perpetual cost). All thermalization figures use the case's own \textbf{plain
HMC, no topological updates} as the standard-practice comparison --- the
instanton move used to build unbiased references in
Section~\ref{sec:appendix-gen} is this pipeline's own fix and is
deliberately excluded here.

\begin{figure}[htbp]
  \centering
  \includegraphics[width=0.95\textwidth]{figures/07_therm_costs.png}
  \caption{\textbf{Diffusion-seed thermalization time vs. two standard-HMC
    costs, full $\beta_f$ range, Wilson-loop observables (plaquette,
    $W(2\times2)$, $W(4\times4)$).} \emph{Left:} raw diffusion-seed
    $t_{\mathrm{therm}}$ (blue) vs. fresh hot-start (red) and cold-start
    (purple) burn-in --- open markers where a fresh chain fails to reach
    tolerance within a 640-trajectory budget. \emph{Right:} the same seed
    curve against the hot-start chain's own steady-state interval
    $2\tau_{\mathrm{int}}$ (gold). Top strip: whether that case's hot-start
    chain's $Q$ tunneled at all (green) or was completely frozen (red).
    \textbf{How derived:} every point comes from running real batched HMC
    (Omelyan integrator, adapted step size) from three starting points ---
    the diffusion seed, a fresh hot start, a fresh cold start --- at matched
    couplings across the generalization study, then measuring each
    observable's trajectory-by-trajectory approach to the exact value.
    \textbf{Significance:} below $\beta_f \approx 8$ (green strip) all
    curves are comparable --- standard HMC needs no help. Above it (red
    strip, the majority of the tested range), both fresh-start burn-in and
    the equilibrated interval grow sharply (burn-in eventually diverging
    outright; the interval reaching 70--90 trajectories at the hardest
    couplings), while the diffusion seed's $t_{\mathrm{therm}}$ stays at
    single digits to low tens of trajectories throughout the entire
    three-decade range --- this is the direct trajectory-cost evidence for
    ``the diffusion seed thermalizes faster than the time it takes standard
    HMC to produce its own next independent sample,'' not just faster than
    burn-in.}
  \label{fig:therm-costs}
\end{figure}

\begin{figure}[htbp]
  \centering
  \includegraphics[width=0.85\textwidth]{figures/08_therm_speedup.png}
  \caption{\textbf{Speedup: burn-in(fresh start)~/~$t_{\mathrm{therm}}$(diffusion
    seed), one pair of bars per case.} Red = vs.\ fresh hot start, purple =
    vs.\ fresh cold start; shaded region = hot-start topology is frozen at
    that coupling; $\geq$ marks a lower bound where the fresh chain never
    reached tolerance within budget (so the true speedup is higher still).
    \textbf{Significance:} turns Figure~\ref{fig:therm-costs} into a single
    number per case --- below freezing a modest 1.3--3$\times$; above it,
    measured lower bounds run from $\sim$11$\times$ up to $>$1000$\times$ at
    couplings where the diffusion seed's raw output is already within
    tolerance before any HMC continuation at all.}
  \label{fig:therm-speedup}
\end{figure}

\subsubsection{Why standard HMC's equilibrium ensemble is wrong, not just slow}

\begin{figure}[htbp]
  \centering
  \includegraphics[width=0.95\textwidth]{figures/09_freezing_zscore.png}
  \caption{\textbf{$\langle Q^2\rangle$ $z$-score vs.\ $\beta_f$: the
    pipeline vs.\ two independent plain-HMC ensembles.} Ensembles built the
    way a practitioner actually would --- one $Q$ sample per independent
    chain, taken after burn-in (repeat samples from a frozen chain carry no
    new information, so this is the correct unit of evidence, not a trick).
    The pipeline (blue) stays inside $|z| \lesssim 2$ across the full
    three-decade range. The hot-start ensemble crosses out of tolerance by
    $\beta_f \approx 6$ and settles at $z \approx 5$--$7$ (a stable bias:
    each frozen chain lands in a different wrong sector, so the ensemble
    mean is systematically off). The cold-start ensemble is worse in a
    different way --- above $\beta_f \approx 9$, essentially every chain
    freezes at exactly $Q=0$, giving zero measured variance and a formally
    enormous $z$ (the honest statement is ``no measured topological
    fluctuation at all,'' not merely ``biased''). \textbf{Significance:}
    demonstrates that the failure above is not a speed problem that more
    trajectories would fix --- it's structural. No feasible trajectory
    budget changes this outcome for standard HMC.}
  \label{fig:freezing-zscore}
\end{figure}

\begin{figure}[htbp]
  \centering
  \includegraphics[width=0.95\textwidth]{figures/10_freezing_pq.png}
  \caption{\textbf{Full $P(Q)$ shape at three couplings spanning ergodic
    $\to$ onset $\to$ frozen.} At $\beta_f = 3.1$ (thousands of tunneling
    events per chain) all three ensembles track the exact shape --- nothing
    to fix here. By $\beta_f = 6.1$ (tunneling already down two orders of
    magnitude) the cold ensemble is visibly over-concentrated at $Q=0$ and
    the hot ensemble is patchy --- early symptoms. At $\beta_f = 219$ (zero
    tunneling events measured in over ten thousand post-burn-in
    chain-trajectories) the breakdown is total: the hot ensemble is
    scattered across sectors with almost no mass where exact actually puts
    it, the cold ensemble is a single zero-width spike at $Q=0$.
    \textbf{Significance:} shows \emph{why} the pipeline avoids this
    failure --- its topological sector is set structurally
    (Section~\ref{sec:appendix-methods}: transported from the coarse
    parent, not reached by letting a chain tunnel), so the one mechanism
    that breaks standard HMC here is simply not in its critical path.}
  \label{fig:freezing-pq}
\end{figure}

\clearpage
\subsubsection{Case studies: low, mid, and high \texorpdfstring{$\beta$}{beta} in detail}

Each case below shows three things for the same physical coupling: how fast
the raw diffusion seed and fresh HMC chains approach the exact value
(relaxation trace), what the seed's distribution looks like against the
case's own plain hot-start HMC \emph{before} any polish, and what it looks
like \emph{after} a short (96-trajectory) HMC continuation --- the same
continuation every pipeline output receives before being reported in
Section~\ref{sec:appendix-gen}.

\paragraph{Low $\beta = 1.49$, $L=32$ (honest baseline).}

\begin{figure}[htbp]
  \centering
  \includegraphics[width=0.95\textwidth]{figures/11_low_relax.png}
  \caption{\textbf{Relaxation trace, $\beta_f = 1.49$.} Rows: plaquette,
    $W(2\times2)$, $W(4\times4)$, $Q^2$; left column zooms the
    early-trajectory approach to exact, right column shows distance from
    exact in error-bar units, log-log, full budget. \textbf{Significance:}
    hot and cold starts both thermalize in single-digit trajectories,
    matching the diffusion seed --- standard HMC is genuinely competitive
    here, confirming the comparison isn't cherry-picked.}
  \label{fig:low-relax}
\end{figure}

\begin{figure}[htbp]
  \centering
  \includegraphics[width=0.95\textwidth]{figures/12_low_before.png}
  \caption{\textbf{Distribution match before any HMC polish, $\beta_f =
    1.49$.} Exact $\langle Q^2\rangle = 28.52$; generated $15.0 \pm 2.0$
    (statistical noise from a modest batch, not bias); the case's own
    hot-start HMC $23.7 \pm 6.2$ --- also fine, since tunneling is still
    fast here.}
  \label{fig:low-before}
\end{figure}

\begin{figure}[htbp]
  \centering
  \includegraphics[width=0.95\textwidth]{figures/13_low_after.png}
  \caption{\textbf{Same case, after the 96-trajectory continuation.}
    Plaquette $z$ moves $-48.3 \to -0.6$ (tightening already-good
    statistics), $\langle Q^2\rangle$ $z$ $-6.8 \to +0.8$.
    \textbf{Significance of this trio (Figures~\ref{fig:low-relax}--\ref{fig:low-after}):}
    at low $\beta$ the standard chain's own burn-in (7 hot / 9 cold
    trajectories) is actually \emph{shorter} than the 96-trajectory polish
    used here --- the diffusion route offers no advantage because none is
    needed. Included for honesty, not to make a case.}
  \label{fig:low-after}
\end{figure}

\paragraph{Mid $\beta = 55.0$, $L=32$ (established freezing).}

\begin{figure}[htbp]
  \centering
  \includegraphics[width=0.95\textwidth]{figures/14_mid_relax.png}
  \caption{\textbf{Relaxation trace, $\beta_f = 55.0$} (case D\_bc14.1464,
    same case as Figures~\ref{fig:mid-before}--\ref{fig:mid-after}). The
    diffusion seed is already inside tolerance at essentially zero
    additional trajectories on the plaquette, $W(2\times2)$, and
    $W(4\times4)$ rows (blue curve sits inside the shaded $|z|\le 2$ band
    from the first trajectory in the right-hand panels); the hot-start
    chain's $Q^2$ (bottom row) plateaus near 10--12, far from the exact
    value $\approx 0.47$, for the entire 640-trajectory budget ---
    established freezing, not just onset.}
  \label{fig:mid-relax}
\end{figure}

\begin{figure}[htbp]
  \centering
  \includegraphics[width=0.95\textwidth]{figures/15_mid_before.png}
  \caption{\textbf{Distribution match before polish, $\beta_f = 55.0$.}
    Exact $\langle Q^2\rangle = 0.474$; generated $0.484 \pm 0.056$
    ($z=+0.18$, matching almost exactly); the case's own hot-start HMC
    $11.19 \pm 1.76$ ($z=+6.09$ --- $23\times$ the true value). The $P(Q)$
    panel shows the mechanism directly: generated mass sits tightly on
    $Q \in \{-2,\ldots,2\}$ matching exact, while the frozen hot ensemble
    is smeared flat across $Q \in \{-7,\ldots,5\}$ with no extra weight
    where the true distribution peaks.}
  \label{fig:mid-before}
\end{figure}

\begin{figure}[htbp]
  \centering
  \includegraphics[width=0.95\textwidth]{figures/16_mid_after.png}
  \caption{\textbf{Same case, after the 96-trajectory continuation.}
    Plaquette $z$ $-2.3 \to -0.3$; $W(2\times2)$ $-0.7 \to +1.6$; $\langle
    Q^2\rangle$ unchanged at $+0.18$ --- the topological sector doesn't
    move in 96 trajectories at this coupling, nor does it need to, since it
    was already correct. \textbf{Significance of this trio
    (Figures~\ref{fig:mid-relax}--\ref{fig:mid-after}):} the raw seed's
    $t_{\mathrm{therm}}$ was already 0 on every Wilson-loop observable, so
    nothing is ``improving'' here so much as staying correct --- while the
    standard chain's cold start alone needs 336 trajectories just to burn
    in, and its hot start is completely frozen (Figure~\ref{fig:freezing-zscore}).}
  \label{fig:mid-after}
\end{figure}

\paragraph{High $\beta = 218.6$, $L=64$ (the money shot).}

\begin{figure}[htbp]
  \centering
  \includegraphics[width=0.95\textwidth]{figures/17_high_relax.png}
  \caption{\textbf{Relaxation trace, $\beta_f = 218.6$, $L=64$} --- same
    case as Figures~\ref{fig:gen-high} and~\ref{fig:gen-volume}. The
    diffusion seed sits inside tolerance from the first trajectory on every
    Wilson observable (blue curve never leaves the shaded band, top three
    rows). The hot-start chain's $Q^2$ (bottom row) plateaus near 30--40
    for the \emph{entire} 640-trajectory budget, nowhere near the exact
    value $\approx 0.47$ --- zero tunneling events, a permanently wrong
    sector.}
  \label{fig:high-relax}
\end{figure}

\begin{figure}[htbp]
  \centering
  \includegraphics[width=0.95\textwidth]{figures/18_high_before.png}
  \caption{\textbf{Distribution match before polish.} Exact $\langle
    Q^2\rangle = 0.474$; generated $0.391 \pm 0.061$ ($z=-1.37$, correct
    sector structure already); the case's own hot-start HMC $52.2 \pm
    22.3$ --- about $110\times$ the true value (the huge error bar, from 16
    independently frozen chains landing in wildly different sectors, is
    itself part of the failure: chain-to-chain scatter masks how wrong the
    mean is).}
  \label{fig:high-before}
\end{figure}

\begin{figure}[htbp]
  \centering
  \includegraphics[width=0.95\textwidth]{figures/19_high_after.png}
  \caption{\textbf{Same case, after the 96-trajectory continuation.} The
    cleanest improvement in this appendix: plaquette $z$ $-7.7 \to +1.4$,
    $W(2\times2)$ $-10.4 \to +1.6$, $W(4\times4)$ $-12.7 \to +1.4$ --- all
    three cross from clear outliers into the passing band. $\langle
    Q^2\rangle$ unchanged at $-1.37$ (already correct; see
    Figure~\ref{fig:high-relax}, bottom row, for the frozen hot-start
    chain's $Q^2$ over the same window it never escapes).
    \textbf{Significance of this trio (Figures~\ref{fig:high-relax}--\ref{fig:high-after}):}
    at the flagship coupling and volume, a 96-trajectory polish (a small
    fraction of any standard-HMC cost measured in
    Section~\ref{sec:appendix-therm}) takes a seed with a real residual
    Wilson-loop bias to comfortably passing, while a fresh hot- or
    cold-start chain never reaches tolerance within the 640-trajectory
    budget at all --- the clearest single illustration that this is not
    ``faster,'' it is ``standard HMC does not arrive at all within any
    practical budget, and the diffusion route does.''}
  \label{fig:high-after}
\end{figure}

\clearpage
% ===================================================================
\subsection{Autocorrelation: fast vs.\ slow modes, over a window of separations}
\label{sec:appendix-autocorr}
% ===================================================================

Everywhere above, ``thermalization'' and ``tunneling'' were described in
terms of a single summary number ($t_{\mathrm{therm}}$, $\tau_{\mathrm{int}}$,
tunneling counts). This section shows the underlying curve those numbers are
extracted from: the normalized autocorrelation function $\Gamma(\delta)$ of
an observable as a function of separation $\delta$ (in HMC trajectories),
computed on the equilibrated window of each chain (post-$t_{\mathrm{therm}}$,
or the second half of the run for a start that never thermalizes), averaged
over chains with SEM error bars. $\Gamma(0) = 1$ by construction; a
well-mixing chain decays toward $\Gamma(\delta) \approx 0$ as $\delta$ grows,
and a flat $\Gamma(\delta)$ that never decays is the direct signature of a
mode the chain cannot move at all within the window tested. \textbf{How
derived:} both figures below reuse
\texttt{diffusion.validate.stats.normalized\_autocorrelation} and the
per-trajectory series already saved by the thermalization benchmark
(\texttt{\{label\}\_series.npz}) --- no new HMC was run for this section,
only new plots (\texttt{diffusion/scripts/11\_autocorrelation.py}) over data
already used elsewhere in this appendix (same three cases as
Figures~\ref{fig:low-relax}--\ref{fig:high-after}: $\beta_f = 1.49$, $55.0$,
and $218.6/L{=}64$).

\begin{figure}[htbp]
  \centering
  \includegraphics[width=0.95\textwidth]{figures/20_autocorr_modes.png}
  \caption{\textbf{Fast modes (Wilson loops, blue) vs.\ slow modes
    (topology, red) side by side, for all three starts (rows) and all three
    couplings (columns).} \textbf{Physics:} the plaquette is a local,
    fast-mixing observable --- it only requires a few nearby link angles to
    change to decorrelate --- while the topological charge is a genuinely
    global mode: changing it requires an instanton-like rearrangement of
    the whole field, which plain HMC (no $Q$-hops) simply does not do at
    appreciable coupling. \textbf{What the panels show:} at $\beta_f = 1.49$
    (left column) both modes decorrelate within about 5 trajectories for
    every start --- weak coupling, nothing frozen. At $\beta_f = 55.0$ and
    $218.6$ (center, right columns), the blue (plaquette) curves still
    decay toward zero for every start, just more slowly at higher $\beta$
    --- ordinary critical slowing down. The red ($Q$) curves, by contrast,
    sit pinned at $\Gamma(\delta) = 1.000$ out to the full 30-trajectory
    window tested, for \textbf{every} start including the diffusion seed's
    own post-generation continuation. \textbf{Significance:} this is not a
    defect in the diffusion route --- once generation is done, the seed is
    handed to the same plain-HMC continuation used everywhere in this
    appendix (no $Q$-hops), so of course its topological charge cannot move
    during that continuation either; nothing here relies on it moving. The
    correctness demonstrated in Section~\ref{sec:appendix-gen} and
    Figures~\ref{fig:freezing-zscore}--\ref{fig:freezing-pq} comes from the
    sector being set correctly at generation time (structural transport,
    Section~\ref{sec:appendix-methods}), not from letting it mix
    afterward. What this figure adds is the direct visual proof that the
    flat-red/decaying-blue pattern is a real, measured property of the
    dynamics --- not an inference from a single $\tau_{\mathrm{int}}$
    number --- and that it is identical in kind for all three starts:
    plain HMC's topological mode simply does not move at these couplings,
    regardless of where the chain began.}
  \label{fig:autocorr-modes}
\end{figure}

\begin{figure}[htbp]
  \centering
  \includegraphics[width=0.95\textwidth]{figures/21_autocorr_fastmode.png}
  \caption{\textbf{Plaquette (fast-mode) autocorrelation only, by start, at
    the same three couplings.} \textbf{Significance:} isolates the
    critical-slowing-down claim from Section~\ref{sec:appendix-therm} on
    its own axis. At $\beta_f = 1.49$ all three starts decay together. At
    $\beta_f = 55.0$ and $218.6$, the diffusion seed's continuation (blue)
    still decays to $\Gamma(\delta) \approx 0$ within about 10 trajectories,
    while the fresh hot-start chain (red) decays far more slowly and
    plateaus at a substantial residual correlation ($\approx 0.1$--$0.4$)
    even after 30 trajectories --- direct, window-resolved evidence for the
    $\tau_{\mathrm{int}}$ growth quoted as a single number in
    Figure~\ref{fig:therm-costs}'s right panel: the fresh chain is still
    meaningfully correlated with itself at a separation where the
    diffusion-seeded continuation has already forgotten its own recent
    history.}
  \label{fig:autocorr-fastmode}
\end{figure}

\clearpage
% ===================================================================
\subsection{Instanton updates vs.\ standard HMC: fixing topological freezing at the source}
\label{sec:appendix-instanton}
% ===================================================================

Sections~\ref{sec:appendix-therm}--\ref{sec:appendix-autocorr} established
\emph{that} plain HMC's topological charge freezes and \emph{why} (a global
tunneling event whose rate collapses with $\beta$). This section isolates
the fix on its own terms, independent of the diffusion pipeline entirely:
\textbf{both ensembles below are pure HMC} --- the same
\texttt{diffusion.lgt.hmc.BatchedHMC} Omelyan integrator, same lattice
($L=32$), same coupling, same step size and trajectory count, same random
hot start --- differing in exactly one variable: whether each step also
attempts the global instanton move
(\texttt{diffusion.lgt.local\_updates.topological\_update}, which proposes
adding a smooth $Q \to Q \pm 1$ configuration to the whole lattice, accepted
with probability $\min(1, e^{-\Delta S})$). No score network, no checkpoint,
no generation of any kind is involved in this section.

\paragraph{How derived.} 32 independent chains per method per coupling, 500
burn-in + 2000 recorded trajectories, $\beta \in \{2, 4, 8, 16, 32, 64, 128,
256\}$ (\texttt{diffusion/scripts/13\_instanton\_vs\_standard\_hmc.py}).
Because within-chain autocorrelation grows severely at high $\beta$ (that is
the whole phenomenon being measured), every mean and error bar below is
computed from the 32 per-chain time-averages --- the only unit that is
rigorously independent regardless of how slowly an individual chain mixes
--- never by pooling all (time $\times$ chain) samples into one estimator,
which would silently overstate precision.

\begin{figure}[htbp]
  \centering
  \includegraphics[width=0.8\textwidth]{figures/22_instanton_tunneling.png}
  \caption{\textbf{Topological-charge tunneling count vs.\ $\beta$, log-log,
    both methods (2000 trajectories $\times$ 32 chains).} Standard HMC:
    54,997 $\to$ 18,662 $\to$ 18 $\to$ \textbf{0} tunnelings from $\beta = 2
    \to 4 \to 8 \to 16$, and stays at exactly zero through $\beta = 256$.
    Instanton HMC: 57,214 at $\beta = 2$, still 943 at $\beta = 256$ --- four
    orders of magnitude more tunneling than standard HMC's zero, at the
    hardest coupling tested. \textbf{Physics:} the collapse from $\beta = 4$
    to $\beta = 8$ is already a $\sim$1000$\times$ drop; continuing that
    trend puts the expected count at $\beta = 16$ well below 1 per 64,000
    chain-trajectories, so landing on exactly 0 is the predicted outcome of
    exponential suppression, not a discontinuity or a counting artifact
    (confirmed directly: at $\beta = 16$ each of the 32 chains sits at a
    \emph{different} frozen integer charge, ranging $-7$ to $+7$, matching
    their random hot start --- each individually has exactly zero variance
    over the full 2000-trajectory window, the signature of real freezing
    rather than a degenerate all-collapse-to-zero bug). \textbf{Significance:}
    this is the controlled version of
    Figure~\ref{fig:freezing-zscore} --- same failure mode, but with the
    diffusion pipeline removed entirely, isolating the instanton move itself
    as the fix.}
  \label{fig:instanton-tunneling}
\end{figure}

\begin{figure}[htbp]
  \centering
  \includegraphics[width=0.8\textwidth]{figures/23_instanton_acceptance.png}
  \caption{\textbf{HMC step acceptance vs.\ instanton-move acceptance,
    vs.\ $\beta$.} The base Omelyan step's acceptance ($\approx$0.96--0.98)
    is statistically identical whether or not the instanton move is enabled
    --- it never touches the leapfrog trajectory, so this is the sanity
    check that adding it doesn't disturb the base sampler. The instanton
    move's own acceptance decays smoothly from 90.6\% ($\beta = 2$) to 1.6\%
    ($\beta = 256$). \textbf{Physics:} the instanton proposal is
    \emph{global} (it adds a fixed smooth configuration to every link at
    once), so its action cost scales as $\Delta S \sim O(\beta/V)$ rather
    than the $O(\beta)$ local-barrier cost that suppresses standard HMC's
    tunneling --- the same $\beta$ range that drives standard HMC's rate to
    exactly zero only drives the instanton move's acceptance down by a
    factor of $\sim$57, not to zero. \textbf{Significance:} this is the
    mechanistic reason Figure~\ref{fig:instanton-tunneling} looks the way it
    does --- a slowly-decaying global move against an
    exponentially-collapsing local one.}
  \label{fig:instanton-acceptance}
\end{figure}

\begin{figure}[htbp]
  \centering
  \includegraphics[width=0.95\textwidth]{figures/24_instanton_traces_standard.png}
  \caption{\textbf{Standard HMC: single representative chain's $Q(t)$, all
    eight couplings stacked.} Direct visual counterpart to
    Figure~\ref{fig:instanton-tunneling}'s aggregate counts. Visibly active
    random walk at $\beta = 2$ and $4$; by $\beta = 8$ the shown chain is
    already flat for its recorded window; from $\beta = 16$ through $256$
    every panel is a dead flat line --- the chain never leaves its starting
    sector for 2000 trajectories. \textbf{Significance:} removes any doubt
    that Figure~\ref{fig:instanton-tunneling}'s zeros are a
    summary-statistic artifact --- the raw trace is visibly, unambiguously
    frozen.}
  \label{fig:instanton-traces-standard}
\end{figure}

\begin{figure}[htbp]
  \centering
  \includegraphics[width=0.95\textwidth]{figures/25_instanton_traces_instanton.png}
  \caption{\textbf{Instanton HMC: the same single-chain view, same eight
    couplings.} Visibly active tunneling at every coupling shown, including
    $\beta = 256$ (24 tunnelings on this specific chain over the recorded
    window, consistent with the pooled 943 across all 32 chains in
    Figure~\ref{fig:instanton-tunneling}). \textbf{Significance:} the
    qualitative contrast with Figure~\ref{fig:instanton-traces-standard} at
    identical couplings is the clearest single picture in this section ---
    same lattice, same $\beta$, same integrator, one move added, and a chain
    that was completely inert is now visibly alive.}
  \label{fig:instanton-traces-instanton}
\end{figure}

\begin{figure}[htbp]
  \centering
  \includegraphics[width=0.95\textwidth]{figures/26_instanton_distributions.png}
  \caption{\textbf{Full observable comparison at $\beta = 64$ ($L=32$):
    standard HMC vs.\ instanton HMC vs.\ exact.} $\langle Q^2 \rangle$:
    standard $22.8 \pm 5.1$ against an exact value of $0.404$ --- wrong by a
    factor of $\sim$56, driven by 32 chains frozen at 32 different,
    mostly-wrong sectors (visible directly in the $P(Q)$ panel, top right:
    standard's mass is smeared across $Q \in \{-4,\ldots,4\}$ with almost
    none at the true peak). Instanton: $0.413 \pm 0.004$ --- correct to
    within 2\%. \textbf{A caveat on the other three panels:} plaquette and
    both Wilson loops show statistically large $z$-scores for \emph{both}
    methods at this $\beta$ ($|z|$ up to $\sim$15), but the underlying
    absolute numbers tell a more precise story --- instanton's mean is
    consistently closer to exact than standard's at every observable (e.g.\
    $W(4\times4)$: standard $0.671$ vs.\ exact $0.882$, 24\% low; instanton
    $0.762$ vs.\ $0.882$, 14\% low), it is just that instanton's per-chain
    scatter is often tighter, which inflates its $z$-score for the same or
    smaller absolute error. The residual bias present in \emph{both} methods
    at this $\beta$ most likely reflects a shared limitation --- the fixed
    500-trajectory burn-in budget probably stops being fully sufficient for
    these observables too as $\beta$ grows, a generic HMC cost independent
    of the instanton move --- rather than something the instanton move
    itself introduces or fully cures. \textbf{Significance:} the topological
    story ($Q$, $\langle Q^2 \rangle$, $P(Q)$) is unambiguous and dramatic;
    the non-topological story is a smaller, genuine, and honestly-reported
    secondary effect --- instanton HMC is consistently the better-matching
    ensemble even there, but by degree, not by a clean pass/fail line the
    way the topological observables show.}
  \label{fig:instanton-distributions}
\end{figure}
